\subsection{10. Chiral Spin Order in Kondo-Heisenberg Systems}
\textbf{OSTI:} \texttt{1412756}\quad\textbf{Authors:} Unknown (OSTI 1412756) (2017)\quad\textbf{Domain:} Condensed Matter Theory / Strongly Correlated Electrons\\
\scorebadge{Coverage}{7}\quad\scorebadge{Agreement}{8}\quad\textbf{Total:} 15/20\par\smallskip
\noindent\textbf{Coverage rationale.} Reproduced the paper's analytical mean-field backbone: zero-T energy functional E\_0(alpha), critical Heisenberg coupling J\_c from C(J\_c)=1 condition, three-regime phase progression (disordered / CSL / AFM), finite-T Ising reduction, and scalar-chirality CSL dome. NOT replicated: rigorous 2D Ising universality (beta=1/8) via field theory or Monte Carlo; lattice-specific triangular/kagome applications; material-specific calculation for Sr2VO3FeAs.\par
\noindent\textbf{Agreement rationale.} J\_c matches analytical Eq. (Jc) to machine precision. Canting angle alpha*(J\_H) shows all three regimes. Quadratic T\_c \textasciitilde{} (J\_H - J\_c)\textasciicircum{}2/J\_K captured near threshold. Phase boundary qualitatively matches paper Fig. 3. Main discrepancy: our mean-field Ising reduction gives weakly first-order transition in some parameter ranges vs paper's claimed continuous 2D Ising transition (beta=1/8) --- acknowledged honestly as MF artifact.\par\smallskip
\noindent\textbf{What reproduced:}
\begin{itemize}[leftmargin=1.4em,itemsep=0.2em,topsep=0.2em]
\item Energy functional E\_0(alpha) (paper Eq. E0)
\item Critical coupling J\_c via C(J\_c)=1 (paper Eq. Jc)
\item Canting angle sin(alpha*) vs J\_H/J\_c
\item Finite-T Ising order parameter m(T) and T\_c(J\_H)
\item Quadratic onset near threshold: T\_c \textasciitilde{} (J\_H-J\_c)\textasciicircum{}2
\item Scalar chirality O\_c \textasciitilde{} s\textasciicircum{}3 sin(alpha) cos\textasciicircum{}2(alpha) CSL dome
\item Z\_2 (time-reversal x parity) symmetry breaking with SU(2) intact
\end{itemize}\noindent\textbf{What missing:}
\begin{itemize}[leftmargin=1.4em,itemsep=0.2em,topsep=0.2em]
\item 2D Ising critical exponent beta=1/8 (requires MC / field theory)
\item Lattice-specific triangular and kagome extensions
\item Material-specific Sr2VO3FeAs numbers
\item Berry curvature / anomalous Hall transport predictions
\item Fluctuation corrections beyond mean-field saddle-point
\end{itemize}\noindent\textbf{Follow-on research questions:}
\begin{enumerate}[leftmargin=1.6em,itemsep=0.2em,topsep=0.2em,label=Q\arabic*.]
\item Run a classical Monte Carlo on the Ising-alpha effective model (on a 2D lattice with coarse-grained coherence volumes) to numerically extract beta and test the claimed 2D Ising universality.
\item Generalize the mean-field ansatz to triangular and kagome Kondo-Heisenberg lattices and compute how the CSL dome shifts --- does kagome frustration widen the CSL window between J\_c and J\_AFM
\item Incorporate itinerant-electron Berry curvature in the CSL phase and predict the anomalous Hall conductivity sigma\_xy(T) below T\_c; compare to measured sigma\_xy in candidate chiral-spin materials.
\item Relax the nested-FS assumption (tilde\_J\_H(Q)=0) and scan filling to map how the CSL dome deforms away from perfect nesting; identify the filling window where CSL survives.
\item Couple the mean-field electronic structure to phonons and ask whether the CSL onset couples to a lattice distortion (magnetoelastic signature) observable in Bragg-peak splitting.
\end{enumerate}

